\chapter{The Stack and the Source}\label{ch:gdb-stack}

You are stopped. The first question is always the same: how did I get here?

\section{Reading a Backtrace}\label{sec:gdb-bt}

\begin{definition}[Call stack]
	The chain of function calls currently in progress. Each entry is a
	\vocab{frame}, holding that call's arguments, its local variables, and the
	address to return to. Frame \texttt{\#0} is where the program is now; the
	numbers increase towards \cmd{main}.
\end{definition}

\begin{keys}[0.26]
	\cmd{backtrace} (\cmd{bt}) & Print the whole stack                            \\
	\cmd{bt 5}                 & The innermost five frames                        \\
	\cmd{bt -5}                & The outermost five                               \\
	\cmd{bt full}              & \dots{} with every frame's local variables       \\
	\cmd{bt -frame-arguments all} & Print full argument values, not summaries     \\
	\cmd{where}                & A synonym for \cmd{bt}                           \\
\end{keys}

\begin{gdbcode}
(gdb) bt
#0  0x00005555555551a9 in parse_line (s=0x0) at parser.c:42
#1  0x0000555555555241 in read_config (path=0x7fffffffe5c1 "app.conf")
    at config.c:18
#2  0x00005555555552c8 in main (argc=2, argv=0x7fffffffe3b8) at main.c:11
\end{gdbcode}

Read it from the bottom. \cmd{main} called \cmd{read\_config} with a filename;
that called \cmd{parse\_line} with \texttt{s=0x0}; and \texttt{0x0} is
\texttt{NULL}, which is the bug. The arguments printed in each frame are
usually enough to locate the fault without looking at a single line of source.

\begin{moral}
	\cmd{bt} is the first command after any crash, and often the last. The frame
	that is \emph{wrong} is rarely frame \texttt{\#0} --- frame \texttt{\#0} is
	merely where the wrongness became fatal. Look for the frame that passed the bad
	value, which is usually one or two up.
\end{moral}

\begin{note}[Frames you cannot read]
	\texttt{??} in a backtrace means no debug information for that frame ---
	a stripped library (\autoref{sec:gdb-symbols}), or a corrupted stack. If
	\emph{every} frame is \texttt{??} and the addresses look implausible, the stack
	itself has been overwritten, and \autoref{sec:gdb-stackoverflow} is the
	relevant section.
\end{note}

\section{Selecting Frames}\label{sec:gdb-frame}

The \vocab{selected frame} is the one that \cmd{print}, \cmd{info locals} and
\cmd{list} refer to. Moving between frames is how you inspect the caller's
variables.

\begin{keys}[0.26]
	\cmd{frame 2} (\cmd{f 2}) & Select frame 2                                     \\
	\cmd{up}                  & One frame outwards, towards \cmd{main}             \\
	\cmd{down}                & One frame inwards                                  \\
	\cmd{up 3}                & Three at once                                      \\
	\cmd{frame}               & Which frame am I in?                               \\
	\cmd{info frame}          & Its addresses, saved registers, return address     \\
\end{keys}

\begin{eg}
\begin{gdbcode}
(gdb) frame 1
#1  0x0000555555555241 in read_config (path=0x7fffffffe5c1 "app.conf")
    at config.c:18
18          parse_line(next_line(f));
(gdb) print f
$1 = (FILE *) 0x0
\end{gdbcode}
	One \cmd{frame 1} and the real bug appears: \cmd{fopen} failed, nobody checked,
	and \cmd{next\_line} was handed a null \texttt{FILE~*}. The crash was two
	frames and one function away from its cause.
\end{eg}

\begin{note}
	Selecting a frame changes nothing about the program --- it only changes which
	frame's names GDB resolves. \cmd{up} and \cmd{down} are safe to use freely, and
	\cmd{frame 0} returns you to the innermost.
\end{note}

\section{Viewing Source}\label{sec:gdb-list}

\begin{keys}[0.26]
	\cmd{list} (\cmd{l})      & Ten lines around the current one                   \\
	\cmd{list}                & Again: the next ten                                \\
	\cmd{list 40,60}          & An explicit range                                  \\
	\cmd{list parse\_line}    & Around a function                                  \\
	\cmd{list parser.c:42}    & Around a line of a file                            \\
	\cmd{list -}              & The ten lines before the last listing              \\
	\cmd{set listsize 20}     & Show more at a time                                \\
	\cmd{directory ../src}    & Where to look for sources GDB cannot find          \\
	\cmd{info source}         & Which file the current frame came from             \\
\end{keys}

\begin{note}
	\texttt{No such file or directory} when listing means the binary records a path
	that no longer exists --- you moved the tree, or built it in a container.
	\cmd{directory} adds a search path, and \cmd{set substitute-path /build /home/konyi/src}
	rewrites a prefix, which is the usual fix for a binary built elsewhere.
\end{note}

\begin{moral}
	If you find yourself typing \cmd{list} more than twice, open the TUI instead
	(\autoref{sec:gdb-tui}). \key{<C-x>} \key{a} turns on a source window that
	follows the program automatically, and the \cmd{list} command becomes
	unnecessary.
\end{moral}

\section{Locals and Arguments}\label{sec:gdb-localargs}

\begin{keys}[0.26]
	\cmd{info locals}      & Every local variable in the selected frame            \\
	\cmd{info args}        & Its arguments                                         \\
	\cmd{info all-registers} & Every register, including floating point            \\
	\cmd{bt full}          & \cmd{info locals} for every frame at once             \\
	\cmd{info functions parse} & Functions whose names match                       \\
	\cmd{info variables}   & Global and static variables                           \\
	\cmd{info scope parse\_line} & What is in scope there, and where each lives     \\
\end{keys}

\begin{gdbcode}
(gdb) info args
s = 0x0
len = 128
(gdb) info locals
i = 0
buf = "\000\000\000..."
p = 0x7fffffffe2a0 "name=value"
\end{gdbcode}

\begin{note}
	\cmd{bt full} is the single most informative command for a crash you cannot
	reproduce: it captures every frame with every local in one block of text, which
	is exactly what to paste into a bug report. Pair it with
	\cmd{set pagination off} so it does not stop every screenful.
\end{note}

\section{Inlined and Optimised Frames}\label{sec:gdb-optimised}

At \cmd{-O2} the stack GDB shows and the stack the machine has are no longer the
same thing, and it is worth knowing how the difference presents itself.

\begin{keys}[0.34]
	\texttt{<optimized out>}        & The variable has no storage at this point. It may be in a register that has since been reused, or eliminated entirely \\
	\texttt{\dots{} (inlined)}      & GDB reconstructed this frame from debug info; no real call happened          \\
	A missing frame                 & A tail call: the compiler reused the caller's frame                          \\
	A line that ``executes twice''  & Instructions from several source lines interleaved                           \\
\end{keys}

\begin{gdbcode}
(gdb) bt
#0  hash (key=<optimized out>) at hash.c:12
#1  0x000055555555529a in lookup (t=0x5555555592a0, key=0x4006f4 "port")
    at table.c:30
#2  table_get (t=0x5555555592a0, key=0x4006f4 "port") at table.c:44
\end{gdbcode}

Note that frame \texttt{\#1} and \texttt{\#2} share an address: \cmd{lookup} was
inlined into \cmd{table\_get}, and GDB is showing you the source-level call that
the machine no longer makes.

\begin{keys}[0.32]
	\cmd{info frame}             & Is this a real frame or a reconstructed one?     \\
	\cmd{set print frame-info source-and-location} & More detail on each frame      \\
	\cmd{p \$rdi}, \cmd{p \$rsi} & Read arguments from registers when names fail    \\
	\cmd{disassemble /s}         & What is actually executing                       \\
\end{keys}

\begin{moral}
	Rebuild at \cmd{-Og} or \cmd{-O0} before fighting an optimised backtrace. If
	the bug then disappears, that is not a dead end --- it is strong evidence of
	undefined behaviour, and \autoref{sec:gdb-companions} will find it far faster
	than reading disassembly would.
\end{moral}
